% Praktikumsanleitung Versuch 1 — Mecanum-Ansteuerung im 2D-Simulator mit ROS 2.
% Bauen: make  (latexmk -pdf). Zwei TikZ-Figuren, keine externen Grafiken nötig.
\documentclass[12pt,a4paper]{scrartcl}
\usepackage[ngerman]{babel}
\usepackage[T1]{fontenc}
\usepackage[utf8]{inputenc}
\usepackage{lmodern}
\usepackage[a4paper,left=2.8cm,right=2.6cm,top=3.0cm,bottom=2.8cm]{geometry}
\usepackage{amsmath,amssymb}
\usepackage{booktabs}
\usepackage{graphicx}
\usepackage[table]{xcolor}
\usepackage{listings}
\usepackage{tikz}
\usetikzlibrary{positioning,calc,arrows.meta,patterns,shapes.geometric}
\usepackage[font=small,skipt=4pt]{caption}
\usepackage{enumitem}
\usepackage{microtype}
\usepackage[colorlinks=true,linkcolor=blue!55!black,urlcolor=blue!55!black,
            citecolor=blue!55!black]{hyperref}
\usepackage{cleveref}

\lstset{language=Python, basicstyle=\ttfamily\small, columns=fullflexible,
        keepspaces=true, tabsize=4, showstringspaces=false, breaklines=true,
        keywordstyle=\bfseries, commentstyle=\itshape\color{black!55},
        stringstyle=\color{black!65}, frame=l, framerule=0.3pt,
        rulecolor=\color{black!30}, xleftmargin=6pt, framexleftmargin=6pt,
        aboveskip=8pt, belowskip=6pt, numbers=none, escapechar=|}

\newcommand{\thema}[1]{\texttt{#1}}
\newcommand{\ordner}[1]{\texttt{#1}}
\newcommand{\kommando}[1]{\texttt{#1}}
\setlist[enumerate,1]{leftmargin=*,itemsep=2pt}
\setlist[itemize,1]{leftmargin=*,itemsep=1pt,topsep=3pt}
\hyphenation{Mechanik Messwerte}

\title{\textbf{Mecanum-Ansteuerung im Simulator}\\[2pt]
  \large Praktikumsversuch 1 — Intelligente Robotik}
% TODO(integrator): Betreuer und Datum eintragen
\author{Betreuung: Robotik-Labor}
\date{Version 0.1 — Wintersemester}

\begin{document}
\maketitle

\begin{abstract}
\noindent
Sie schreiben die \emph{Ansteuerung der vier Mecanum-Räder} eines mobilen Roboters.
Dazu übersetzen Sie eine gewünschte Körpergeschwindigkeit in vier Radgeschwindigkeiten
in \(\mathrm{rad/s}\) und veröffentlichen sie auf ein ROS-Thema. Ein kleiner
2D-Simulator (Pygame, rund 1400 Zeilen Python) fährt Ihren Code in einer von drei
Umgebungen und liefert Odometrie, LIDAR und eine globale Position zurück.
Der Simulator ist absichtlich winzig: Sie können die Physik in einer halben Stunde
lesen und sind damit in der Lage, ihn für eigene Fragen zu verändern.
\end{abstract}

\tableofcontents

\section{Versuchsziel und Vorfragen}

\subsection{Ziel}
Nach dem Versuch können Sie
\begin{itemize}
  \item die Kinematik eines Mecanum-Fahrwerks aus der Rolleranordnung herleiten
        und die Vorzeichen begründen (statt sie zu kopieren),
  \item den Unterschied zwischen inverse und forward Kinematik sowie deren
        pseudo-inverse Auflösung erklären,
  \item einen ROS-2-Knoten schreiben, der Themen liest (\thema{odom}, \thema{scan},
        \thema{gps}) und veröffentlicht (\thema{cmd\_vel}, \thema{wheel\_speeds}),
  \item die Drift einer Dead-Reckening-Odometrie messen und einordnen.
\end{itemize}
Gearbeitet wird allein pro Person: ein Robotername pro Teilnehmer, ein Knoten pro
Roboter. Der Bewerter misst ausschließlich über die Themen — wie Sie implementieren,
steht Ihnen frei.

\subsection{Vorfragen (vor dem Praktikumsstermin rechnen)}
Verwenden Sie die Simulatorwerte aus \cref{tab:parameter}:
\(l_x = 0{,}14\,\mathrm{m}\), \(l_y = 0{,}13\,\mathrm{m}\), \(r = 0{,}05\,\mathrm{m}\),
\(a = l_x + l_y = 0{,}27\,\mathrm{m}\), maximale Radgeschwindigkeit
\(w_{\max} = 12\,\mathrm{rad/s}\). Lösungen: \cref{sec:loesungen}.

\begin{enumerate}
  \item \textbf{Vorwärts.} Wie schnell dreht jedes der vier Räder bei
        \(v_x = 0{,}4\,\mathrm{m/s}\), \(v_y = \omega = 0\)?
  \item \textbf{Seitwärts.} Gesucht \(v_y = +0{,}25\,\mathrm{m/s}\) (nach links).
        Berechnen Sie die vier Raddrehzahlen und erklären Sie, warum zwei Räder
        negativ sind, obwohl der Roboter nicht dreht.
  \item \textbf{Drehen auf der Stelle.} Für \(\omega = +1{,}2\,\mathrm{rad/s}\)
        (gegen den Uhrzeigersinn): welche Vorzeichen haben VL und HL — und warum
        dieselben?
  \item \textbf{Diagonalfahrt.} \(v_x = v_y = 0{,}3\,\mathrm{m/s}\). Ein Rad steht
        still. Welches? Welche Konsequenz hat das für die höchstfahrbare
        Diagonalgeschwindigkeit, wenn alle Räder mit \(w_{\max}\) begrenzt sind?
  \item \textbf{Rückwärtsrechnen.} Die Räder messen
        \([w_\text{FL}, w_\text{FR}, w_\text{RL}, w_\text{RR}] =
        [8, 0, 0, 8]\,\mathrm{rad/s}\). Bestimmen Sie \(v_x\), \(v_y\), \(\omega\)
        und geben Sie Fahrtrichtung und Betrag an.
\end{enumerate}

\section{Das Mecanum-Prinzip}
\subsection{Rolleranordnung}
Ein Mecanum-Rad ist ein normales Rad, dessen Mantel aus schräg (45°) montierten
Rollen besteht. Es kann damit (i) wie ein normales Rad angetrieben in seine
Wirkrichtung \(d\) rollen und (ii) senkrecht dazu, \(f\), frei gleiten.
In der Draufsicht zeigen die vier Wirkrichtungen nach vorn-innen; es entsteht ein X —
die \emph{X-Anordnung} (Abbildung~\ref{fig:mecanum}).

\begin{figure}[ht]
  \centering
  % TikZ-Figur: Mecanum-Fahrwerk von oben.
% Die vier Wirkrichtungen d_i sind aus der Konvention in CONTRACT 5 abgeleitet
% (r*w_FL = vx - vy  =>  d_FL = (1,-1) usf.) — der Text leitet das erst her.
\begin{tikzpicture}[
  scale=0.92, every node/.style={font=\small},
  wheel/.style={draw=black!75, line width=0.7pt, fill=black!10,
                minimum width=1.35cm, minimum height=0.58cm, inner sep=0pt},
  drive/.style={-{Stealth[length=2mm]}, draw=blue!60!black, line width=0.9pt},
  free/.style={-{Stealth[length=2mm]}, draw=black!40, dashed, line width=0.7pt},
  roller/.style={draw=black!55, line width=0.5pt},
  axis/.style={-{Stealth[length=2.6mm]}, line width=1pt},
  dim/.style={-{Stealth[length=1.5mm]}-{Stealth[length=1.5mm]}, draw=black!55,
              line width=0.5pt, font=\footnotesize},
  box/.style={draw=black!60, rounded corners=2pt, fill=yellow!12, inner sep=3pt,
              font=\footnotesize, align=left},
]
\def\lx{1.75}   % halbe Fahrzeuglaenge in Bild Einheiten (kein Massstab!)
\def\ly{1.25}   % halbe Fahrzeugbreite

% Karosserie
\draw[line width=0.8pt, rounded corners=3pt, fill=black!4]
  (-\lx-0.45,-\ly-0.42) rectangle (\lx+0.45,\ly+0.42);
\node[font=\footnotesize\itshape, text=black!55] at (0.55,0) {Basisplatte};

% Rad: Position und Winkel der Wirkrichtung in Grad (Labels stehen separat)
\newcommand{\mwheel}[3]{%
  \begin{scope}[shift={(#1,#2)}]
    \node[wheel] at (0,0) {};
    % Rollen: schraffierte Streifen senkrecht auf der Freigleit-Richtung
    \begin{scope}[rotate={#3+90}]
      \draw[roller] (-0.42,-0.26) -- (-0.18,0.26);
      \draw[roller] (-0.06,-0.26) -- (0.18,0.26);
      \draw[roller] ( 0.30,-0.26) -- ( 0.54,0.26);
    \end{scope}
    \draw[free,  rotate={#3+90}] (-0.05,0) -- (0.95,0)
      node[right=0pt, font=\scriptsize, text=black!55] {$f$};
    \draw[drive, rotate=#3] (-0.05,0) -- (1.05,0)
      node[right=0pt, font=\scriptsize, text=blue!60!black] {$d$};
  \end{scope}}

\mwheel{ \lx}{ \ly}{-45}   % VL: vorne links
\mwheel{ \lx}{-\ly}{ 45}   % VR: vorne rechts
\mwheel{-\lx}{ \ly}{ 45}   % HL
\mwheel{-\lx}{-\ly}{-45}   % HR

\fill[black] (0,0) circle (1.4pt);
\node[anchor=north east, font=\footnotesize] at (0,-0.1) {Schwerpunkt};

% Welt-/Körperachsen
\draw[axis] (0,0) -- (2.9,0)   node[right, font=\footnotesize] {$x$ (vorn)};
\draw[axis] (0,0) -- (0,2.15)  node[above, font=\footnotesize] {$y$ (links)};
\draw[axis, draw=black!65] (0.95,0) arc (0:74:0.95) node[right, font=\footnotesize] {$\theta$ (CCW $+$)};

% Spurmass und Radstand
\draw[dim] (-\lx,-\ly-0.78) -- (-\lx,-\ly-0.18);
\draw[dim] ( \lx,-\ly-0.78) -- ( \lx,-\ly-0.18);
\draw[dim] (-\lx,-\ly-0.62) -- ( \lx,-\ly-0.62) node[midway, above=1pt] {$2\,l_x$ (Radstand)};
\draw[dim] (\lx+0.72,-\ly) -- (\lx+0.72,\ly);
\node[rotate=90, font=\footnotesize, text=black!70] at (\lx+0.98,0) {$2\,l_y$ (Spurweite)};

% Radlabels
\node[font=\footnotesize, anchor=south] at ( \lx, \ly+0.42) {$w_\text{FL}$ VL};
\node[font=\footnotesize, anchor=north] at ( \lx,-\ly-0.42) {$w_\text{FR}$ VR};
\node[font=\footnotesize, anchor=south] at (-\lx, \ly+0.42) {$w_\text{RL}$ HL};
\node[font=\footnotesize, anchor=north] at (-\lx,-\ly-0.42) {$w_\text{RR}$ HR};

% Randnotiz mit der Konvention
\node[box, anchor=north west] at (0.2,-3.15) {%
  Rand-und-Vorn-Regel: $x$ nach vorn, $y$ nach \emph{links}, $\theta$ gegen den Uhrzeigersinn\\
  Rollachsen-Pfeile $d$ ergeben in der Draufsicht ein X ($X$-Anordnung).\\
  $f$ = Freigleit-Richtung (dreht das Fahrzeug auf der Stelle, bleibt unvermittelt).\\
  Radreihenfolge in allen Themen: $[w_\text{FL}, w_\text{FR}, w_\text{RL}, w_\text{RR}]
   = [\text{VL}, \text{VR}, \text{HL}, \text{HR}]$};
\end{tikzpicture}

  \caption{Mecanum-Fahrwerk von oben. Blaue Pfeile \(d\) = angetriebene
    Wirkrichtung, gestrichelte Pfeile \(f\) = Freigleit-Richtung, schraffierte
    Streifen = Rollen. Achsen und Vorzeichen sind exakt die des Simulators.}
  \label{fig:mecanum}
\end{figure}

\subsection{Geschwindigkeitszerlegung an einem Rad}
Der Schwerpunkt habe die Körpergeschwindigkeit \(\boldsymbol{\xi} =
(v_x, v_y, \omega)^{\mathsf T}\); das Rad \(i\) sitze an der Stelle
\(\mathbf p_i = (\pm l_x, \pm l_y)\). Sein Kontaktpunkt bewegt sich mit
\begin{equation}
  \mathbf v_i \;=\; \begin{pmatrix} v_x \\ v_y\end{pmatrix}
  + \omega \begin{pmatrix} -p_{i,y} \\ p_{i,x}\end{pmatrix},
  \qquad\text{z.\,B. für VL:}\quad
  \mathbf v_\text{FL} = \begin{pmatrix} v_x - l_y\omega \\ v_y + l_x\omega\end{pmatrix}.
\end{equation}
Das Rad liefert nur in Wirkrichtung Drehzahl, in Freigleit-Richtung überträgt es
nichts. Mit den Wirkrichtungen aus \cref{fig:mecanum},
\(d_\text{FL} = (1,-1)\), \(d_\text{FR} = (1,1)\), \(d_\text{RL} = (1,1)\),
\(d_\text{RR} = (1,-1)\) (Projektion ohne Normierung — der Faktor \(\sqrt2\) steckt in
der effektiven Radradius-Kalibrierung), folgt \(r\,w_i = \mathbf v_i \cdot d_i\):
\begin{align}
  r\,w_\text{FL} &= (v_x - l_y\omega) - (v_y + l_x\omega) = v_x - v_y - a\,\omega
   \label{eq:fl}\\
  r\,w_\text{FR} &= (v_x + l_y\omega) + (v_y - l_x\omega) = v_x + v_y + a\,\omega
   \label{eq:fr}\\
  r\,w_\text{RL} &= (v_x - l_y\omega) + (v_y + l_x\omega) = v_x + v_y - a\,\omega
   \label{eq:rl}\\
  r\,w_\text{RR} &= (v_x + l_y\omega) - (v_y - l_x\omega) = v_x - v_y + a\,\omega
   \label{eq:rr}
\end{align}
mit dem Gier-Hebelarm \(a = l_x + l_y\). Das ist die \textbf{inverse Kinematik} — Ihre
Aufgabe in T1. Als Matrix (Radreihenfolge \([\)FL, FR, RL, RR\(]\)):
\begin{equation}
  \mathbf w = J\,\boldsymbol{\xi},\qquad
  J = \frac{1}{r}\begin{pmatrix}
    1 & -1 & -a\\ 1 & 1 & a\\ 1 & 1 & -a\\ 1 & -1 & a
  \end{pmatrix}.
  \label{eq:J}
\end{equation}

\subsection{Vorwärtskinematik als Pseudo-Inverse}
Gesucht ist \(\boldsymbol\xi\) aus gemessenen Raddrehzahlen, also \(\mathbf w = J
\boldsymbol\xi\) mit vier Gleichungen und drei Unbekannten. Da \(J\) Spaltenrank 3 hat,
ist die Moore-Penrose-Lösung \(^+\)die Pseudo-Inverse
\(J^{+} = (J^{\mathsf T}J)^{-1}J^{\mathsf T}\). Es ist
\(J^{\mathsf T}J = \tfrac{1}{r^2}\operatorname{diag}(4,4,4a^2)\) — die drei
Freiheitsgrade sind also \emph{entkoppelt}, und es ergibt sich
\begin{equation}
  J^{+} = \frac{r}{4}\begin{pmatrix}
    1 & 1 & 1 & 1\\ -1 & 1 & 1 & -1\\
    -\tfrac{1}{a} & \tfrac{1}{a} & -\tfrac{1}{a} & \tfrac{1}{a}
  \end{pmatrix}
  \;\Rightarrow\;
  \begin{aligned}
    v_x    &= \tfrac{r}{4}\,(w_\text{FL} + w_\text{FR} + w_\text{RL} + w_\text{RR})\\
    v_y    &= \tfrac{r}{4}\,(-w_\text{FL} + w_\text{FR} + w_\text{RL} - w_\text{RR})\\
    \omega &= \tfrac{r}{4a}\,(-w_\text{FL} + w_\text{FR} - w_\text{RL} + w_\text{RR})
  \end{aligned}
  \label{eq:fk}
\end{equation}
Der Simulator benutzt \cref{eq:fk} für die Odometrie: Was die Räder tatsächlich
machen, entscheidet, wo der Roboter ankommt — nicht das, was Sie wollten. Deshalb
wirkt sich ein falsches Vorzeichen in T1 sofort als Fehlfahrt aus.

\begin{quote}\small
\textbf{Vorzeichen-Box (Merksatz).}\ \newline
\(x\) nach vorn, \(y\) nach \textbf{links}, \(\theta\) gegen den Uhrzeigersinn positiv.
Radreihenfolge in \emph{allen} Themen: \([\)FL, FR, RL, RR\(^] = [\text{VL}, \text{VR},
\text{HL}, \text{HR}]\). Drehzahl in \(\mathrm{rad/s}\), Körpergeschwindigkeit in
\(\mathrm{m/s}\) bzw. \(\mathrm{rad/s}\).
\end{quote}

\section{Der Simulator}
\subsection{Architektur}
\begin{figure}[ht]
  \centering
  % TikZ-Figur: Datenfluss im Simulator — von der ASCII-Karte bis zum Radkommando.
\begin{tikzpicture}[
  every node/.style={font=\small},
  blk/.style={draw=black!70, line width=0.6pt, rounded corners=2pt, fill=black!5,
              align=center, inner sep=4pt, minimum height=8.6mm},
  ros/.style={blk, fill=blue!7, draw=blue!50!black},
  stu/.style={blk, fill=green!9, draw=green!45!black, very thick},
  fl/.style={-{Stealth[length=2mm]}, line width=0.7pt, draw=black!70},
  flw/.style={-{Stealth[length=2.4mm]}, line width=1.1pt, draw=green!40!black},
  lbl/.style={font=\scriptsize\ttfamily, inner sep=1.5pt, fill=white, fill opacity=0.9},
]
\node[blk] (world) {worlds/\texttt{*.txt} \\ (ASCII-Grid)};
\node[blk, right=10mm of world] (eng) {\textbf{engine.SimEngine}\\[1pt]
  \scriptsize Takt, Sensorraten, Watchdog};
\node[blk, right=10mm of eng] (phys) {physics.Chassis\\ \scriptsize Wahrheit, Kollision};
\node[blk, below=10mm of eng] (sen) {sensors\\ \scriptsize Odometry, Lidar, GpsSensor};
\node[ros, below=10mm of sen] (bus) {ros\_bridge.RclpyBus \quad$\vert$\quad stub.StubBus};
\node[blk, below=7mm of phys, xshift=4mm] (gui) {render.py\\ \scriptsize Pygame-HUD};
\node[stu, below=11mm of bus] (io) {robot\_io.RobotIO\\ \scriptsize Messwerte rein, Räder raus};
\node[stu, right=10mm of io] (ctrl) {student/controller\\\_template.py\\
  \scriptsize inverse\_kinematics, mission};

\draw[fl] (world) -- (eng) node[lbl, midway] {Wände, Starts, Ziel};
\draw[fl] (eng) -- (phys) node[lbl, midway] {\texttt{step(dt)}};
\draw[fl] (phys.west) -- ++(-0.4,0) |- (sen.east)
  node[lbl, near end, below] {Pose, Rad-Istwerte};
\draw[fl] (sen) -- (bus) node[lbl, midway] {Messungen};
\draw[fl] (eng.south west) |- (bus.north west);
\draw[fl] (eng.east) |- node[lbl, pos=0.3, above, sloped] {read only} (gui.west);
\draw[fl] (bus) -- (io) node[lbl, midway] {ROS-Themen};
\draw[flw] (io) -- (ctrl);
\draw[flw] (ctrl.east) -- ++(0.7,0) |- node[lbl, pos=0.62, right]
  {\texttt{<robot>/wheel\_speeds}} (bus.east);

\node[anchor=north west, font=\scriptsize\ttfamily, text=black!70]
  at ($(sen.north west)+(0.25,0.26)$) {odom 50\,Hz \quad scan 20\,Hz \quad gps 5\,Hz};
\node[anchor=north west, font=\scriptsize\itshape, text=black!60, align=left]
  at ($(bus.south west)+(-0.15,-1.05)$)
  {Ohne ROS übernimmt \texttt{stub.StubBus} alle Kanten — derselbe Studentencode.\\
   \texttt{cmd\_vel} (Twist) läuft gegenläufig: Bewerter oder Teleop $\to$ Sim.};
\end{tikzpicture}

  \caption{Datenfluss. Der geschlossene Regelkreis läuft über \thema{wheel\_speeds}:
    Nur Ihre vier Radzahlen bestimmen die Pose. Der Renderer ist rein lesend.}
  \label{fig:architektur}
\end{figure}

Der Simulator existiert in zwei Spielarten, am selben Code:
\emph{mit} ROS 2 (\ordner{ros\_bridge.py} → echter DDS-Bus) und \emph{ohne} ROS
(\ordner{stub.py} → Bus im selben Prozess). Für Sie ist der Studentencode identisch.

\subsection{Themen und Typen}
Frames: \thema{map} (Welt), \thema{odom} (Odometrie-Ursprung), \thema{base\_link}
(Roboter), \thema{laser}. Roboternamen tragen ihre Themen als Namensraum.

\begin{table}[ht]\centering\small
\caption{ROS-Schnittstelle (Pfad \thema{<robot>} = Robotername).}
\label{tab:themen}
\begin{tabular}{@{}p{4.3cm}p{4.9cm}p{5.2cm}@{}}
\toprule
Thema & Typ (ROS 2) & Bedeutung \\
\midrule
\thema{<robot>/cmd\_vel} & \thema{geometry\_msgs/msg/Twist}
  & Soll-Körpergeschwindigkeit; \(v_y>0\) heißt \emph{links} \\
\thema{<robot>/wheel\_speeds} & \thema{std\_msgs/msg/Float64MultiArray} (4)
  & Ihre vier Radzahlen in \(\mathrm{rad/s}\) — der Übungsteil \\
\thema{<robot>/odom} & \thema{nav\_msgs/msg/Odometry}
  & Dead Reckoning aus \emph{gemessenen} Raddrehzahlen, verrauscht \\
\thema{<robot>/scan} & \thema{sensor\_msgs/msg/LaserScan}
  & 360 Strahlen über \(2\pi\); Strahl 0 zeigt in Fahrtrichtung, wachsender Index
  gegen den Uhrzeigersinn; kein Treffer \(=\) \thema{range\_max} \\
\thema{<robot>/gps} & \thema{geometry\_msgs/msg/PoseStamped}
  & globale Position (UWB/MoCap-Charakter), verrauscht \\
\thema{<robot>/truth} & \thema{geometry\_msgs/msg/PoseStamped}
  & exakte Pose — nur, wenn \thema{debug\_truth} eingeschaltet ist \\
\thema{<robot>/mission\_state} & \thema{std\_msgs/msg/String}
  & \thema{idle}, \thema{running}, \thema{done}, \thema{failed:<grund>} \\
\thema{/sim/robots} & \thema{std\_msgs/msg/String} (JSON)
  & alle Roboter mit Farbe, Marker, Motorvariante, Modus \\
\thema{/sim/world} & \thema{std\_msgs/msg/String} (JSON)
  & Weltname, Größe, Ziel, Startposes \% TODO(integrator): in CONTRACT 3 ergänzen \\
\thema{/sim/task} & \thema{std\_msgs/msg/String}
  & aktiver Auftrag des Bewerters \\
\thema{/sim/spawn\_robot} & Service, siehe \cref{sec:service}
  & Roboter hinzufügen (eindeutiger Name Pflicht) \\
\thema{/sim/despawn\_robot} & Service, siehe \cref{sec:service}
  & Roboter entfernen \\
\thema{/sim/reset} & \thema{std\_srvs/srv/Trigger} & alle Roboter zurück an den Start \\
\thema{/clock} & \thema{rosgraph\_msgs/msg/Clock} & Simulationszeit (\thema{use\_sim\_\_time}) \\
\bottomrule
\end{tabular}
\end{table}

\subsection{Was verrauscht, was exakt ist}
\begin{table}[ht]\centering\small
\caption{Parameter aus \ordner{config/default.json} (Auszug).}
\label{tab:parameter}
\begin{tabular}{@{}llp{7.3cm}@{}}
\toprule
Größe & Wert & Bemerkung \\
\midrule
Physikrate & 50 Hz & feste Schrittweite, damit Läufe reproduzierbar sind \\
Odometrie & 50 Hz & Rauschen \(\sigma_{\text{wheel}}=0{,}04\,\mathrm{rad/s}\) auf
  jeder Radmessung \(\Rightarrow\) integrierte Drift; \(\sigma_{xy}=1{,}5\,\mathrm{mm}\),
  \(\sigma_\theta=1{,}2\,\mathrm{mrad}\) auf dem Ausgang \\
LIDAR & 20 Hz & 360 Strahlen, \(0\ldots8\,\mathrm{m}\), \(\sigma=15\,\mathrm{mm}\);
  Roboter sehen sich in Versuch 1 nicht \\
Globale Position & 5 Hz & \(\sigma_{xy}=6\,\mathrm{cm}\), \(\sigma_\theta=0{,}03\,\mathrm{rad}\) \\
Motor & \(\tau=0{,}06\,\mathrm{s}\) & Trägheit 1. Ordnung,
  \(|\dot w|\le 40\,\mathrm{rad/s^2}\), \(|w|\le 12\,\mathrm{rad/s}\) \\
Kommando-Wachdog & 0{,}35 s & ohne neue Meldungrollen die Räder aus \\
Kollision & Kreis \(r=0{,}21\,\mathrm{m}\) & Position wird aus der Wand geschoben,
  Berührungen zählen (\thema{contacts}) \\
\bottomrule
\end{tabular}
\end{table}

Exakt (deterministisch, kein Rauschen) sind nur die \emph{Wahrheit} in
\ordner{physics.py} und — falls eingeschaltet — das Thema \thema{<robot>/truth}.
Schlupf wird nicht modelliert: vier Radzahlen ergeben eindeutig
\((v_x,v_y,\omega)\) nach \cref{eq:fk}.

\subsection{pass-through-Modus und wheels-Modus}
Damit der erste Tag nicht an der Kinematik scheitert, startet jeder Roboter im
\textbf{pass-through}-Modus: Kommt \thema{cmd\_vel} und \emph{nie} eine
\thema{wheel\_speeds}-Meldung, rechnet die Simulation die Twist selbst in vier Räder —
der Roboter fährt, als wäre Ihre Lösung schon richtig. Senden Sie die erste
\thema{wheel\_speeds}-Meldung, kippt der Roboter für immer in den
\textbf{wheels}-Modus: jetzt gilt ausschließlich, was Sie rechnen (Modus steht in
\thema{/sim/robots} und im HUD). Ein kaputtes Vorzeichen ist damit sofort sichtbar —
und genau so gemeint.

\section{Inbetriebnahme}
\subsection{Installation}
\begin{lstlisting}[language=bash]
sudo apt install python3-pygame          # Ubuntu, reicht fuer den Stub-Modus
git clone <repo>                          # folgt im naechsten Schritt des Praktikums
cd mecanum-lab
./lab run --world maze --robot alice      # Sim + Ihr Knoten, EIN Prozess, kein ROS
\end{lstlisting}
Mit ROS 2 (getestet mit Kilted/Jazzy, Ubuntu 24.04):
\begin{lstlisting}[language=bash]
source /opt/ros/$ROS_DISTRO/setup.bash
./lab sim --world production --robots alice     # Sim solo, veroeffentlicht ROS-Themen
\end{lstlisting}
\subsection{Ausprobieren}
\begin{lstlisting}[language=bash]
./lab teleop --robot alice               # Twist-Fernsteuerung auf cmd_vel
./lab robots                             # wer faehrt, in welchem Modus?
ros2 topic list
ros2 topic echo /alice/odom --once
ros2 topic hz   /alice/scan              # ~20 Hz
ros2 topic echo /alice/wheel_speeds      # Ihre vier Radzahlen
\end{lstlisting}
\subsection{Roboter per Service anlegen}\label{sec:service}
Ein Roboter pro Teilnehmer, Name eindeutig (\thema{[a-z][a-z0-9\_]\{1,23\}}).
\textbf{Achtung ROS 2:} anders als in ROS 1 gibt es \emph{kein} gebautes
\thema{std\_srvs/SetString}. Der Simulator bietet deshalb zwei Wege, die Aufrufer
gleichermaßen funktionieren:
\begin{enumerate}
  \item \textbf{Eigenes Interface} (empfohlen, eine einmalige Übung in
        \texttt{.srv}-Definitionen):
\begin{lstlisting}
# interfaces/mecanum_lab_interfaces/srv/SpawnRobot.srv
string name        # eindeutig
string variant     # "" = automatisch (stock|slow|fast|agile)
---
bool success
string message
uint8 index
string color
string marker
float64 x
float64 y
float64 theta
\end{lstlisting}
\begin{lstlisting}[language=bash]
colcon build --packages-select mecanum_lab_interfaces
source install/setup.bash
ros2 service call /sim/spawn_robot mecanum_lab_interfaces/srv/SpawnRobot \
  "{name: alice, variant: ''}"
\end{lstlisting}
  \item \textbf{JSON-Fallback ohne Interface-Bau}: Handshake über Themen,
        \thema{/sim/spawn\_robot/request} (\thema{std\_msgs/String}, JSON) →
       Antwort auf \thema{/sim/spawn\_robot/result}, Timeout 2\,s.
\end{enumerate}
Beide Wege kapseln \kommando{./lab spawn --name alice} und
\kommando{./lab despawn --name alice}; der Bewerter benutzt ebenfalls diesen Weg.
\subsection{Stolperfallen}
\begin{itemize}
  \item \textbf{Kein DISPLAY} (SSH, CI): \kommando{SDL\_VIDEODRIVER=dummy} oder
        \kommando{./lab sim --headless}.
  \item \textbf{ROS nicht gerodet}: \kommando{ros2} fehlt → alles läuft automatisch
        über den Stub-Bus; \thema{ros2 topic list} zeigt dann natürlich nichts.
  \item \textbf{use\_sim\_\_time}: Knoten mit Wanduhrzeit statt \thema{/clock}
        vergleichen Zeitstempel falsch → \kommando{--ros-args -p use\_sim\_time:=true}.
  \item \textbf{Name bereits vergeben}: \thema{success=false},
        \thema{message="name vergeben"} — erst \kommando{./lab despawn}, dann neu.
  \item \textbf{Namensraum}: Themen hängen am Roboternamen, \thema{/alice/odom} ≠
        \thema{/bob/odom}. Wer keinen Roboternamen übergibt, wartet ins Leere.
  \item \textbf{Funkstille}: ohne Meldung innerhalb 0{,}35\,s rollen die Räder aus —
        Ihr Knoten muss mit \(\ge 10\,\mathrm{Hz}\) senden.
\end{itemize}

\section{Aufgaben}\label{sec:aufgaben}
Alle Schwellen stehen in \ordner{config/tasks.json} und sind dort änderbar; der
Bewerter misst nur über Themen. Ausgeführt wird in der Produktionsumgebung,
\kommando{./lab grade --robot alice --task alle} (Punkte pro Auftrag in
\cref{tab:aufgaben}).

\begin{table}[ht]\centering\small
\caption{Aufträge, Punkte, Erfolgsbedingungen.}
\label{tab:aufgaben}
\begin{tabular}{@{}p{1.35cm}p{6.7cm}p{2.4cm}p{4.2cm}@{}}
\toprule
ID & Inhalt & Grenze & Punkte \\
\midrule
\thema{kinematik} & T1: \thema{cmd\_vel} \(\rightarrow\) vier Räder. Drei Phasen zu je
  3\,s: \(v_x=0{,}3\); \(v_y=0{,}3\); \(\omega=0{,}6\).
  & \begin{tabular}{@{}l@{}}\(\Delta x>0{,}35\,\mathrm{m}\)\\ \(\Delta y>0{,}35\,\mathrm{m}\)\\
  \(\Delta\theta>0{,}9\,\mathrm{rad}\)\\ je \(|\text{Rest}|<0{,}12\,\mathrm{m}\)\end{tabular}
  & 30 \\
\thema{quadrat} & T2: Quadratfahrt, vier Seiten \(1{,}0\,\mathrm{m}\), \(90^\circ\)
  CCW, nur mit Odometrie. & Abschluss \(<0{,}20\,\mathrm{m}\), \(\Delta\theta<15^\circ\),
  Weg \(3{,}2\ldots9{,}0\,\mathrm{m}\), 0 Berührungen, \(<90\,\mathrm{s}\) & 30 \\
\thema{korridor} & T3: Ziel der Welt anfahren, LIDAR verhindert jede Wandberührung.
  & Zielfehler \(<0{,}25\,\mathrm{m}\), 0 Berührungen, seitlicher Abstand
  \(>0{,}16\,\mathrm{m}\) bei Geradaus (\(|\omega|<0{,}25\)) & 30 \\
\thema{gps\_anfahrt} & T4 (Bonus): mit GPS-Rückführung zur eigenen Startpose.
  & Fehler \(<0{,}25\,\mathrm{m}\), 0 Berührungen, \(<90\,\mathrm{s}\) & 10 \\
\bottomrule
\end{tabular}
\end{table}

\subsection*{T1 — Inverse Kinematik (30 Punkte)}
Sie ergänzen in \ordner{student/controller\_template.py} genau eine Funktion. Der
Rest des Knotens (Themen, Takt, Wachdog) ist vorbereitet:
\begin{lstlisting}
def inverse_kinematics(vx, vy, omega):
    """[VL, VR, HL, HR] in rad/s — Konvention aus Abschnitt 2."""
    a = LX + LY                      # TODO: Hebelarm fuer Gier
    return [(vx - vy - a * omega) / R,   # TODO: VL
            (vx + vy + a * omega) / R,   # TODO: VR
            (vx + vy - a * omega) / R,   # TODO: HL
            (vx - vy + a * omega) / R]   # TODO: RR
\end{lstlisting}
\textbf{Teilschritte.} (1) Stub-Modus starten, Teleop benutzen, im HUD die
Räder beobachten. (2) Ein Vorzeichen nach dem anderen abprüfen: nur \(v_x\), dann nur
\(v_y\), dann nur \(\omega\). (3) \kommando{./lab grade --task kinematik} laufen
lassen. \textbf{Lernziel:} Vorzeichen aus der Rollergeometrie verstehen, nicht raten.

Der Simulator selbst rechnet (Auszug aus \ordner{mecanum\_lab/physics.py}, Ihre
Rechnung muss auf dieselben Werte kommen):
\begin{lstlisting}
def inverse_kinematics(g, vx, vy, omega) -> list:
    """Kuerpergeschwindigkeit -> vier Solldrehzahlen [VL, VR, HL, HR] in rad/s."""
    a = g.arm
    return [(vx - vy - a * omega) / g.r, (vx + vy + a * omega) / g.r,
            (vx + vy - a * omega) / g.r, (vx - vy + a * omega) / g.r]

def forward_kinematics(g, wheels) -> tuple:
    """Vier Radgeschwindigkeiten -> (vx, vy, omega); exakte Umkehrung der obigen Zeilen."""
    fl, fr, rl, rr = wheels
    a = g.arm
    return (g.r / 4 * (fl + fr + rl + rr),
            g.r / 4 * (-fl + fr + rl - rr),
            g.r / (4 * a) * (-fl + fr - rl + rr))
\end{lstlisting}

\subsection*{T2 — Quadratfahrt mit Odometrie (30 Punkte)}
Erweitern Sie \thema{mission(rob, task)}: vier Seiten \(1{,}0\,\mathrm{m}\), jeweils
\(90^\circ\) gegen den Uhrzeigersinn, Rückführung über \thema{rob.odom()}.
Melden Sie \thema{done}, wenn Sie wieder an der Startpose sind. Typischer Weg:
Geradeaus-P-Drehregelung auf den Achsenfehler, Totband für die Drehung,
Zwischenziele in der Odometrie gespeichert. Messen Sie zusätzlich Ihren
Abschlussfehler und kommentieren Sie, woher er kommt (Drift, Trägheit, Winkel
darstellung).

\subsection*{T3 — Ziel anfahren, der LIDAR schützt (30 Punkte)}
Ziel: \thema{goal} der Welt (steckt als JSON in \thema{/sim/world}). Richtung und
Strecke aus der Odometrie, Sicherheit aus \thema{rob.scan()}: seitlicher Abstand
\(>0{,}16\,\mathrm{m}\), \emph{keine} Wandberührung (\thema{contacts}=0). Sinnvolle
Bausteine: Vokablen aus Strahlmittelwerten (nicht einzelne Strahlen!),
Geschwindigkeitsreduktion mit dem vorderen Mindestabstand, Not-Stopp.

\subsection*{T4 — Bonus: GPS statt Odometrie (10 Punkte)}
Rückweg zur eigenen Startpose, aber geregelt auf \thema{rob.gps()}. Beobachten und
kommentieren Sie: Totzone eines P-Reglers, Messrauschen als Zittern, 5\,Hz gegen
50\,Hz Odometrie.

\section{Abgabe und Bewertung}
\begin{itemize}
  \item Abgegeben wird \ordner{student/<name>\_controller.py} in Ihrem Branch/Ordner;
        Robotername = Dateiname = Teilnehmerkürzel.
  \item Der Autograder ist maßgeblich:
        \kommando{./lab grade --robot <name> --task alle --json ergebnis.json}.
        Ein Lauf pro Abgabe, Werte aus \cref{tab:aufgaben}.
  \item Dazu eine kurze README (\(\le 20\) Zeilen): Vorgehen, bekannte Probleme,
        was Sie anders machen würden.
  \item Selbsttest vor der Abgabe: \kommando{./lab run --robot test --controller
        student/<name>\_controller.py --headless} und
        \kommando{python3 -m pytest tests -q}.
\end{itemize}
\subsection*{Mussertext für die Commit-Meldung}
\begin{lstlisting}[language=bash, frame=none, xleftmargin=0pt]
T2: Quadratfahrt mit Odometrie (1.0 m, 90 Grad CCW)

- inverse_kinematics auf Konvention Abschnitt 2 korrigiert (vy vorzeichenrichtig)
- Drehrueckfuehrung ueber odom.theta, Totband 0.03 rad, P-Gewinn 1.6
- bekannt: Abschlussfehler ~0.15 m durch Drift; GPS (T4) waere die naechste Stufe
\end{lstlisting}
\subsection*{Bewertungsraster}
\begin{table}[ht]\centering\small
\begin{tabular}{@{}p{7.6cm}p{2.6cm}p{4.2cm}@{}}
\toprule
Kriterium & Punkte & geprüft durch \\
\midrule
T1 Kinematik (drei Phasen je 10 P) & 30 & Autograder \\
T2 Quadratfahrt (Abschluss 20, Fahrqualität 10) & 30 & Autograder \\
T3 Ziel ohne Berührung (Zielfehler 20, Abstand 10) & 30 & Autograder \\
T4 Bonus GPS & 10 & Autograder \\
 kurzes Resümee im README & 0–5 & Betreuung \\
\bottomrule
\end{tabular}
\end{table}
\subsection*{Hilfsmittel und Selbstständigkeit}
Abschreiben von Lösungsteilen (auch aus generierten Antworten anderer Personen oder
Sprachmodelle) führt zum Rücktritt des Versuchs; diskutieren dürfen Sie alles.
Hilfsmittel: Skript, diese Anleitung, \ordner{docs/}, Ihre eigenen Notizen. Die
Musterlösung \ordner{student/solution.py} ist im Repo und wird \emph{erst nach
Abgabeschluss} freigegeben \% TODO(integrator): Freigabezeitpunkt festlegen.
Die vier Aufträge sind hintereinander zu bearbeiten; T1 ist Voraussetzung für alles
andere.

\section{Ausblick}
\begin{itemize}
  \item \textbf{Versuch 2 — Odometrie-Fusion}: Odometrie (50 Hz, Drift) mit GPS
        (5 Hz, verrauscht) zu einem Schätzer kombinieren; Kompass fehlt bewusst.
  \item \textbf{Versuch 3 — Pfadfolgen}: Rundkurs-Welt, Follow-the-Corridor mit
        LIDAR, Geschwindigkeitsbegrenzung durch Kurvenradius.
  \item \textbf{Versuch 4 — Mehrroboter}: mehrere Roboter, eigene Namensräume,
        Kollisionen zwischen Robotern (dafür muss \ordner{sensors.py} erweitert
        werden — genau die Art Änderung, die der kleine Simulator erlaubt).
\end{itemize}

\appendix
\section{Themen-, Typen- und API-Übersicht}
\subsection*{Nachrichtenaufbau (Auszug)}
\begin{lstlisting}
# /alice/wheel_speeds : std_msgs/msg/Float64MultiArray
msg.data = [w_FL, w_FR, w_RL, w_RR]      # rad/s, Reihenfolge ist Vertrag

# /alice/odom : nav_msgs/msg/Odometry
odom.pose.pose.position.{x,y}            # Weltkarte, Framework odom
odom.pose.pose.orientation.z, .w         # Quaternion um z (theta)
odom.twist.twist.linear.{x,y}            # Koerpferam, y nach links
odom.twist.twist.angular.z               # rad/s, CCW positiv

# /alice/scan : sensor_msgs/msg/LaserScan
scan.angle_min = 0.0                     # Strahl 0 = Fahrtrichtung
scan.angle_increment = 2*pi/360          # wachsender Index: CCW
scan.ranges[i]                           # m; kein Treffer = range_max (8.0)

# /alice/gps : geometry_msgs/msg/PoseStamped   (Framework map)
# /sim/world : std_msgs/msg/String, JSON: name, size, goal, spawns, walls
\end{lstlisting}
\subsection*{Studentische API (\texorpdfstring{\ordner{robot\_io.py}}{robot_io.py})}
\begin{lstlisting}
rob = RobotIO("alice")           # ./lab run --robot alice
rob.cmd_vel()                    # Twist oder None
rob.odom(); rob.gps(); rob.scan()
rob.age("odom")                  # Alter der letzten Messung in s
rob.send_wheels([w1, w2, w3, w4])
rob.publish_cmd_vel(vx, vy, omega)
rob.set_mission_state("done")
rob.task(); rob.world(); rob.robots()
serve(module)                    # Runner hinter der einen Zeile im Template
\end{lstlisting}
\textbf{Vertrag mit \thema{serve}:} Ist \thema{/sim/task} leer oder
\thema{kinematik}, wird jede \thema{cmd\_vel} mit \thema{inverse\_kinematics} in
Räder übersetzt. Bei einem anderen Auftrag wird genau einmal \thema{mission(rob,
task)} aufgerufen und der Zustand gemeldet. \% TODO(integrator): finalen Wortlaut des
Templates prüfen

\section{Formelsammlung}
\begin{align}
  \text{IK: } r\begin{pmatrix}w_\text{FL}\\w_\text{FR}\\w_\text{RL}\\w_\text{RR}\end{pmatrix}
  &= \begin{pmatrix}v_x - v_y - a\omega\\ v_x + v_y + a\omega\\
     v_x + v_y - a\omega\\ v_x - v_y + a\omega\end{pmatrix},
   && a = l_x + l_y \\
  \text{FK: } v_x = \tfrac{r}{4}\textstyle\sum_i w_i,\quad
   v_y = \tfrac{r}{4}(-w_\text{FL}+w_\text{FR}+w_\text{RL}-w_\text{RR}),\quad
   \omega = \tfrac{r}{4a}(-w_\text{FL}+w_\text{FR}-w_\text{RL}+w_\text{RR}) \\
  \text{Grenze: } v_{\max} = r\,w_{\max} = 0{,}6\,\mathrm{m/s},\qquad
   \omega_{\max} = \frac{r\,w_{\max}}{a} = 2{,}22\,\mathrm{rad/s}
   \quad(\text{jeweils für } v_y \text{ bzw. reine Drehung}) \\
  \text{Winkel normieren: } \theta \in (-\pi,\pi],\quad
   \Delta\theta = \mathrm{wrap}(\theta_\text{Ziel} - \theta_\text{ist})
\end{align}

\section{Lösungen der Vorfragen}\label{sec:loesungen}
\begin{enumerate}
  \item Alle vier gleich: \(w = v_x/r = 0{,}4/0{,}05 = 8\,\mathrm{rad/s}\).
        \(a\) spielt keine Rolle, \(\omega=0\).
  \item \(w_\text{FL} = -v_y/r = -5\), \(w_\text{FR} = +5\), \(w_\text{RL} = +5\),
        \(w_\text{RR} = -5\,\mathrm{rad/s}\). Die Räder müssen \emph{gegeneinander}
        arbeiten, weil jede Rolle nur längs \(d\) Kraft überträgt: die linke
        Radseite wird nach hinten, die rechte nach vorn geschoben, die Freigleit-
        richtungen lassen die Längsbewegung des Fahrzeugs zu. Kein Giermoment,
        weil sich die Hebel aufheben (\(-,+ ,+,-\) mit Vorzeichen aus \cref{eq:fl,eq:fr,eq:rl,eq:rr}).
  \item \(w = \mp a\omega/r = \mp 0{,}27\cdot1{,}2/0{,}05 = \mp 6{,}48\,\mathrm{rad/s}\),
        konkret \([-6{,}48, +6{,}48, -6{,}48, +6{,}48]\). VL und HL haben dasselbe
        Vorzeichen, weil bei reiner Drehung das Vorzeichen nur von \(x_i\) (Radpaar
        vorne/hinten) und dem Drehsinn abhängt — nicht von links/rechts; die
        Seitwärtsanteile sind hier Null.
  \item \(w_\text{FL} = w_\text{RR} = (0{,}3-0{,}3)/0{,}05 = 0\),
        \(w_\text{FR} = w_\text{RL} = 12\,\mathrm{rad/s}\). Stillstand bei VL/RR ist
        korrekt. Bei \(w_{\max}=12\) ist die Diagonale gerade grenzwertig:
        \((v_x - v_y) = 0\), \((v_x+v_y)=0{,}6 = r\,w_{\max}\) — diagonal ist bei
        \(v_x = v_y = 0{,}3\,\mathrm{m/s}\) Schluss, obwohl längs \(0{,}6\) ginge.
  \item \(v_x = \tfrac{0{,}05}{4}(8+0+0+8) = 0{,}2\,\mathrm{m/s}\),
        \(v_y = \tfrac{0{,}05}{4}(-8+0+0-8) = -0{,}2\,\mathrm{m/s}\) (nach rechts!),
        \(\omega = \tfrac{0{,}05}{4\cdot0{,}27}(-8+0-0+8) = 0\). Betrag
        \(\sqrt{2}\cdot0{,}2 = 0{,}283\,\mathrm{m/s}\), Richtung \(45^\circ\)
        vorn-rechts. Als Soll-Kommando wäre das \((v_x, v_y) = (0{,}2, -0{,}2)\).
\end{enumerate}

\end{document}
